% ----------------------
\subsection{Section 77 (Between circa 4 and circa 20 March 1832)}
% ----------------------
	\label{DC77}
	
	% -----
	\subsubsection{Historical Background}
	% -----
		\label{DC77-HB}
		
		% --
		\paragraph{JSP}
		% --
		
			``As JS continued his revision of the New Testament in February and March 1832, he reached the book of Revelation with its abundance of symbolic language. ``About the first of march, in connection with the translation of the scriptures,'' a later JS history explains, ``I received the following explanation of the Revelations of Saint John.'' Given that JS was in Kirtland, Ohio, and was not working on the New Testament revisions between 29 February and 4 March 1832, this document was likely written sometime between 4 March and 20 March, when another revelation told JS and Sidney Rigdon (who was serving as JS's scribe) to ``omit the translation for the present time'' so that they could travel to Missouri. By the night of 24--25 March 1832, when the pair was attacked by a group of men in Hiram, Ohio, JS and Rigdon were working on the eleventh chapter of the book of Revelation, the last chapter mentioned in the explanation.
			
			``Because he was inscribing the New Testament revision, Rigdon probably served as the original scribe for the explanation, but Jesse Gause could have been the scribe instead. The earliest surviving copy is an undated one made by John Whitmer in Revelation Book 1, where it is identified only as ``Revelation Explained.'' Whitmer, who was residing in Missouri at the time, probably made the copy sometime after April 1832, when JS likely brought a copy of the original document to Missouri along with copies of revelations dictated in March.''
		
	% -----
	\subsubsection{Versions}
	% -----
		\label{DC77-V}
		
		See the \href{https://drive.google.com/file/d/1goAvNz8jS4R8slYfMG_db55Ty2z_vmgK/view?usp=sharing}{sections comparison} document.
	
		\begin{compactdesc}
			\item[JSP] \hyperref[EMS-1832-JS-CIR-4-MAR-1832]{Between circa 4 and circa 20 March 1832}
			\item[BC] ---
			\item[DC 1835] ---
			\item[TS] 5:14 (1 August 1844), 595--596
			\item[DC 1844] ---
			\item[MS] 14:9 (24 April 1852), 132--133
		\end{compactdesc}

	% -----
	\subsubsection{Section 77:1} % *DC77-1 
	% -----
		\label{DC77-1}

		\paragraph{Q.}

			What is the sea of glass spoken of by John, 4th chapter, and 6th verse of the 	Revelation?

		\paragraph{A.}
		
			 It is the earth, in its sanctified, immortal, and eternal state.

		% \begin{framed}
		% \scriptsize
			% \begin{compactdesc}
				% \item[JSP] 
				% % \item[BC] 
				% % \item[]
			% \end{compactdesc}
		% \end{framed}

	% -----
	\subsubsection{Section 77:2} % *DC77-2 
	% -----
		\label{DC77-2}

		\paragraph{Q.}

			What are we to understand by the four beasts, spoken of in the same verse?

		\paragraph{A.}
		
			They are figurative expressions, used by the Revelator, John, in describing heaven, the paradise of God, the happiness of man, and of beasts, and of creeping things, and of the fowls of the air; that which is spiritual being in the likeness of that which is temporal; and that which is temporal in the likeness of that which is spiritual; the spirit of man in the likeness of his person, as also the spirit of the beast, and every other creature which God has created.

		% \begin{framed}
		% \scriptsize
			% \begin{compactdesc}
				% \item[JSP] 
				% % \item[BC] 
				% % \item[]
			% \end{compactdesc}
		% \end{framed}

	% -----
	\subsubsection{Section 77:3} % *DC77-3 
	% -----
		\label{DC77-3}

		\paragraph{Q.}

			Are the four beasts limited to individual beasts, or do they represent classes or orders?

		\paragraph{A.}
		
			They are limited to four individual beasts, which were shown to John, to represent the glory of the classes of beings in their destined order or sphere of creation, in the enjoyment of their eternal felicity.

		% \begin{framed}
		% \scriptsize
			% \begin{compactdesc}
				% \item[JSP] 
				% % \item[BC] 
				% % \item[]
			% \end{compactdesc}
		% \end{framed}

	% -----
	\subsubsection{Section 77:4} % *DC77-4 
	% -----
		\label{DC77-4}

		\paragraph{Q.}

			What are we to understand by the eyes and wings, which the beasts had?

		\paragraph{A.}
		
			Their eyes are a representation of light and knowledge, that is, they are full of knowledge; and their wings are a representation of power, to move, to act, etc.

		% \begin{framed}
		% \scriptsize
			% \begin{compactdesc}
				% \item[JSP] 
				% % \item[BC] 
				% % \item[]
			% \end{compactdesc}
		% \end{framed}

	% -----
	\subsubsection{Section 77:5} % *DC77-5 
	% -----
		\label{DC77-5}

		\paragraph{Q.}

			What are we to understand by the four and twenty elders, spoken of by John?

		\paragraph{A.}
		
			We are to understand that these elders whom John saw, were elders who had been faithful in the work of the ministry and were dead; who belonged to the seven churches, and were then in the paradise of God.

		% \begin{framed}
		% \scriptsize
			% \begin{compactdesc}
				% \item[JSP] 
				% % \item[BC] 
				% % \item[]
			% \end{compactdesc}
		% \end{framed}

	% -----
	\subsubsection{Section 77:6} % *DC77-6 
	% -----
		\label{DC77-6}

		\paragraph{Q.}

			What are we to understand by the book which John saw, which was sealed on the back with seven seals?

		\paragraph{A.}
		
			We are to understand that it contains the revealed will, mysteries, and the works of God; the hidden things of his economy%
				\footnote{%
					% \label{}%
					For ``economy,'' see DC 76, \hyperref[DC76-1-economy]{this note}.
					}
			concerning this earth during the seven thousand years of its continuance, or its temporal existence.

		% \begin{framed}
		% \scriptsize
			% \begin{compactdesc}
				% \item[JSP] 
				% % \item[BC] 
				% % \item[]
			% \end{compactdesc}
		% \end{framed}

	% -----
	\subsubsection{Section 77:7} % *DC77-7 
	% -----
		\label{DC77-7}

		\paragraph{Q.}

			What are we to understand by the seven seals with which it was sealed?

		\paragraph{A.}
		
			We are to understand that the first seal contains the things of the first thousand years, and the second also of the second thousand years, and so on until the seventh.

		% \begin{framed}
		% \scriptsize
			% \begin{compactdesc}
				% \item[JSP] 
				% % \item[BC] 
				% % \item[]
			% \end{compactdesc}
		% \end{framed}

	% -----
	\subsubsection{Section 77:8} % *DC77-8 
	% -----
		\label{DC77-8}

		\paragraph{Q.}

			What are we to understand by the four angels, spoken of in the 7th chapter and 1st verse of Revelation?

		\paragraph{A.}
		
			We are to understand that they are four angels sent forth from God, to whom is given power over the four parts of the earth, to save life and to destroy; these are they who have the everlasting gospel to commit to every nation, kindred, tongue, and people; having power to shut up the heavens, to seal up unto life, or to cast down to the regions of darkness.

		% \begin{framed}
		% \scriptsize
			% \begin{compactdesc}
				% \item[JSP] 
				% % \item[BC] 
				% % \item[]
			% \end{compactdesc}
		% \end{framed}

	% -----
	\subsubsection{Section 77:9} % *DC77-9 
	% -----
		\label{DC77-9}

		\paragraph{Q.}
		
			What are we to understand by the angel ascending from the east, Revelation 7th chapter and 2nd verse?

		\paragraph{A.}
		
			We are to understand that the angel ascending from the east is he to whom is given the seal of the living God over the twelve tribes of Israel; wherefore, he crieth unto the four angels having the everlasting gospel, saying: Hurt not the earth, neither the sea, nor the trees, till we have sealed the servants of our God in their foreheads.  And, if you will receive it, this is Elias which was to come to gather together the tribes of Israel and restore all things.

		% \begin{framed}
		% \scriptsize
			% \begin{compactdesc}
				% \item[JSP] 
				% % \item[BC] 
				% % \item[]
			% \end{compactdesc}
		% \end{framed}

	% -----
	\subsubsection{Section 77:10} % *DC77-10 
	% -----
		\label{DC77-10}

		\paragraph{Q.}

			What time are the things spoken of in this chapter to be accomplished?

		\paragraph{A.}
		
			They are to be accomplished in the sixth thousand years, or the opening of the sixth seal.

		% \begin{framed}
		% \scriptsize
			% \begin{compactdesc}
				% \item[JSP] 
				% % \item[BC] 
				% % \item[]
			% \end{compactdesc}
		% \end{framed}

	% -----
	\subsubsection{Section 77:11} % *DC77-11 
	% -----
		\label{DC77-11}

		\paragraph{Q.}

			What are we to understand by sealing the one hundred and forty-four thousand, out of all the tribes of Israel—twelve thousand out of every tribe?

		\paragraph{A.}
		
			We are to understand that those who are sealed are high priests,%
				\footnote{%
					% \label{}%
					The initial ordinations for high priests had happened less than a year before this, I think.
					}
			ordained unto the holy order of God, to administer the everlasting gospel; for they are they who are ordained out of every nation, kindred, tongue, and people, by the angels to whom is given power over the nations of the earth, to bring as many as will come to the church of the Firstborn.%
				\footnote{%
					% \label{}%
					About a month before this, JS received the vision of \hyperref[DC76]{DC 76}, which mentions ``the church of the Firstborn.''
					}

		% \begin{framed}
		% \scriptsize
			% \begin{compactdesc}
				% \item[JSP] 
				% % \item[BC] 
				% % \item[]
			% \end{compactdesc}
		% \end{framed}

	% -----
	\subsubsection{Section 77:12} % *DC77-12 
	% -----
		\label{DC77-12}

		\paragraph{Q.}

			What are we to understand by the sounding of the trumpets, mentioned in the 8th chapter of Revelation?

		\paragraph{A.}
		
			We are to understand that as God made the world in six days, and on the seventh day he finished his work, and sanctified it, and also formed man out of the dust of the earth, even so, in the beginning of the seventh thousand years will the Lord God sanctify the earth, and complete the salvation of man, and judge all things, and shall redeem all things, except that which he hath not put into his power, when he shall have sealed all things, unto the end of all things; and the sounding of the trumpets of the seven angels are the preparing and finishing of his work, in the beginning of the seventh thousand years—the preparing of the way before the time of his coming.

		% \begin{framed}
		% \scriptsize
			% \begin{compactdesc}
				% \item[JSP] 
				% % \item[BC] 
				% % \item[]
			% \end{compactdesc}
		% \end{framed}

	% -----
	\subsubsection{Section 77:13} % *DC77-13 
	% -----
		\label{DC77-13}

		\paragraph{Q.}
		
			When are the things to be accomplished, which are written in the 9th chapter of Revelation?

		\paragraph{A.}
		
			They are to be accomplished after the opening of the seventh seal, before the coming of Christ.

		% \begin{framed}
		% \scriptsize
			% \begin{compactdesc}
				% \item[JSP] 
				% % \item[BC] 
				% % \item[]
			% \end{compactdesc}
		% \end{framed}

	% -----
	\subsubsection{Section 77:14} % *DC77-14 
	% -----
		\label{DC77-14}

		\paragraph{Q.}

			What are we to understand by the little book which was eaten by John, as mentioned in the 10th chapter of Revelation?

		\paragraph{A.}
		
			We are to understand that it was a mission, and an ordinance, for him to gather the tribes of Israel; behold, this is Elias, who, as it is written, must come and restore all things.

		% \begin{framed}
		% \scriptsize
			% \begin{compactdesc}
				% \item[JSP] 
				% % \item[BC] 
				% % \item[]
			% \end{compactdesc}
		% \end{framed}

	% -----
	\subsubsection{Section 77:15} % *DC77-15 
	% -----
		\label{DC77-15}

		\paragraph{Q.}
		
			What is to be understood by the two witnesses, in the eleventh chapter of Revelation?

		\paragraph{A.}
		
			They are two prophets that are to be raised up to the Jewish nation in the last days, at the time of the restoration, and to prophesy to the Jews after they are gathered and have built the city of Jerusalem in the land of their fathers.

		% \begin{framed}
		% \scriptsize
			% \begin{compactdesc}
				% \item[JSP] 
				% % \item[BC] 
				% % \item[]
			% \end{compactdesc}
		% \end{framed}